\SGANextSectionTarget{sga3:xxiii:section:5}\SGASection{5}{Corollaries of the Fundamental Theorem}

The most important is the following.

\SGARefTarget{sga3:XXIII:5.1}
\SGAStatement{Corollary 5.1}{Let \(S\) be a nonempty scheme, let
\(G,G'\) be pinned \(S\)-groups, and let
\[
  h:R(G')\xrightarrow{\sim}R(G)
\]
be an isomorphism of pinned root data. There is a unique isomorphism
of pinned \(S\)-groups \(f:G\xrightarrow{\sim}G'\) such that
\(R(f)=h\).}

This also follows from \SGARefLink{sga3:XXIII:3.5.1}{Theorem 3.5.1}, because its relations can be
written using only \(R(G)\). It also follows from the elementary part
of the proof of \SGARefLink{sga3:xxiii:theorem:4-1}{Theorem 4.1}, without the exceptional isogenies of
\SGARefLink{sga3:XXIII:4.1.7}{4.1.7} and \SGARefLink{sga3:XXIII:4.1.8}{4.1.8}.

\SGARefTarget{sga3:XXIII:5.2}
\SGAStatement{Corollary 5.2 (Uniqueness theorem)}{Let \(S\) be a
scheme and let \(G,G'\) be two splittable \(S\)-groups. If they have
the same type, then they are isomorphic.}

\SGARefTarget{sga3:XXIII:5.3}
\SGAStatement{Corollary 5.3}{Let \(S\) be a scheme and \(G,G'\) two
splittable \(S\)-groups. The following are equivalent:
\begin{enumerate}
  \item[\SGAInlineRefTarget{sga3:XXIII:5.3:item:i}(i)] \(G\) and \(G'\) are isomorphic;
  \item[\SGAInlineRefTarget{sga3:XXIII:5.3:item:ii}(ii)] they are locally isomorphic for the fpqc topology;
  \item[\SGAInlineRefTarget{sga3:XXIII:5.3:item:iii}(iii)] there is \(s\in S\) such that \(G_s\) and \(G'_s\)
    have the same type.
\end{enumerate}}

Clearly \SGARefLink{sga3:XXIII:5.3:item:i}{(i)} implies \SGARefLink{sga3:XXIII:5.3:item:ii}{(ii)}, which implies \SGARefLink{sga3:XXIII:5.3:item:iii}{(iii)}. Conversely, \SGARefLink{sga3:XXIII:5.3:item:iii}{(iii)}
gives
\[
  R(G)=R(G_s)=R(G'_s)=R(G'),
\]
so \SGARefLink{sga3:XXIII:5.2}{Corollary 5.2} applies.

\SGARefTarget{sga3:XXIII:5.4}
\SGAStatement{Corollary 5.4 (Uniqueness of Chevalley schemes)}{Let
\(G,G'\) be reductive groups over \(\mathbf Z\), each possessing a
split maximal torus. The following are equivalent:
\begin{enumerate}
  \item[(i)] \(G\) and \(G'\) are isomorphic;
  \item[(ii)] there is \(s\in\operatorname{Spec}\mathbf Z\) such
    that \(G_s\) and \(G'_s\) have the same type;
  \item[(iii)] \(G_{\mathbf C}\) and \(G'_{\mathbf C}\) have the same
    type.
\end{enumerate}}

Indeed, \(G\) and \(G'\) are splittable by \SGARefLink{sga3:xxii:prop:2.2}{Exposé XXII, 2.2}.
\NDE{In fact every \(\mathbf Z\)-torus is split: every
\(\mathbf Z\)-torus is isotrivial, and every étale cover of
\(\operatorname{Spec}\mathbf Z\) is trivial.}

\SGARefTarget{sga3:XXIII:5.5}
\SGAStatement{Corollary 5.5 (Existence of outer automorphisms)}{Let
\(S\) be a scheme, \(G\) a pinned \(S\)-group, and \(h\) an
automorphism of the pinned root datum \(R(G)\). There is a unique
automorphism \(u\) of \(G\) preserving the pinning and satisfying
\(R(u)=h\).}

\SGARefTarget{sga3:XXIII:5.5bis}
\SGAStatement{Corollary 5.5 bis}{Let \(S\) be a scheme, \(G\) a
split reductive \(S\)-group, and \(R^+\) a system of positive roots.
For every simple root \(\alpha\), choose an isomorphism of vector
groups
\[
  p_\alpha:\mathbf G_{a,S}\xrightarrow{\sim}U_\alpha.
\]
Let \(h\) be an automorphism of \(M\) permuting the positive roots
and the corresponding coroots: if \(\alpha\in R^+\), then
\[
  h(\alpha)\in R^+,\qquad h^\vee(\alpha^*)=h(\alpha)^*.
\]
There is a unique automorphism \(u\) of \(G\) inducing \(D_S(h)\) on
\(T\) and satisfying \(u\circ p_\alpha=p_{h(\alpha)}\) for every
simple root \(\alpha\).}

\SGARefTarget{sga3:XXIII:5.6}
\SGAStatement{Corollary 5.6}{Let \(G,G'\) be reductive
\(S\)-groups. The following are equivalent:
\begin{enumerate}
  \item[\SGAInlineRefTarget{sga3:XXIII:5.6:item:i}(i)] \(G,G'\) are locally isomorphic for the fpqc topology;
  \item[(i bis)] they are locally isomorphic for the étale topology;
  \item[(ii)] \(G_s,G'_s\) are isomorphic for every \(s\in S\);
  \item[\SGAInlineRefTarget{sga3:XXIII:5.6:item:iii}(iii)] the functions \(s\mapsto\operatorname{type}(G_s)\)
    and \(s\mapsto\operatorname{type}(G'_s)\) are equal.
\end{enumerate}}

The implication (i bis)\(\Rightarrow\)\SGARefLink{sga3:XXIII:5.6:item:i}{(i)} is immediate;
(i)\(\Rightarrow\)(ii) follows from the principle of finite
extension (EGA IV\(_3\), 9.1.4); and
(ii)\(\Rightarrow\)(iii) is immediate. For
\SGARefLink{sga3:XXIII:5.6:item:iii}{(iii)}\(\Rightarrow\)(i bis), one may suppose \(G,G'\) splittable
(\SGARefLink{sga3:xxii:cor:2.3}{Exposé XXII, 2.3}), and then apply \SGARefLink{sga3:XXIII:5.3}{Corollary 5.3}.

\SGARefTarget{sga3:XXIII:5.7}
\SGAStatement{Corollary 5.7}{Let \(S\) be a scheme, \(G\) a
reductive \(S\)-group, and \(G'\) an affine smooth \(S\)-group with
connected fibers. If \(s\in S\) and \(G_s\simeq G'_s\), then there
is an étale morphism \(S'\to S\), covering \(s\), such that
\(G_{S'}\simeq G'_{S'}\).}

By Exposés XIX, \SGARefLink{sga3:XXIII:2.5}{2.5} and XXII, \SGARefLink{sga3:xxii:prop:2.1}{2.1}, one may suppose \(G,G'\)
reductive and splittable, and \SGARefLink{sga3:XXIII:5.3}{Corollary 5.3} applies.

\SGARefTarget{sga3:XXIII:5.8}
\SGAStatement{Corollary 5.8}{Let \(k\) be a field and \(G,G'\)
reductive \(k\)-groups. The following are equivalent:
\begin{enumerate}
  \item[(i)] \(G,G'\) have the same type;
  \item[(ii)] \(G\otimes_k\bar k\) and
    \(G'\otimes_k\bar k\) are isomorphic;
  \item[(iii)] for some finite separable extension \(K/k\),
    \(G\otimes_kK\) and \(G'\otimes_kK\) are isomorphic.
\end{enumerate}}

\SGARefTarget{sga3:XXIII:5.9}
\SGAStatement{Corollary 5.9}{Let \(S\) be a nonempty scheme and
\(R\) a root datum. The following are equivalent:
\begin{enumerate}
  \item[\SGAInlineRefTarget{sga3:XXIII:5.9:item:i}(i)] there is a pinned \(S\)-group of type \(R\);
  \item[(ii)] there is an \(S\)-group of type \(R\);
  \item[\SGAInlineRefTarget{sga3:XXIII:5.9:item:iii}(iii)] fpqc-locally, there is a reductive \(S\)-group of
    type \(R\).
\end{enumerate}}

Only \SGARefLink{sga3:XXIII:5.9:item:iii}{(iii)}\(\Rightarrow\)\SGARefLink{sga3:XXIII:5.9:item:i}{(i)} needs proof. Suppose, for simplicity,
that there are a faithfully flat quasi-compact morphism \(S'\to S\)
and a reductive \(S'\)-group \(G'\) of type \(R\). We may suppose
\(G'\) splittable and choose a pinning \(E'\), writing
\(R=R(G',E')\). The two pullbacks of \((G',E')\) to
\[
  S''=S'\times_SS'
\]
are pinned groups \((G''_1,E''_1)\), \((G''_2,E''_2)\), with
canonical isomorphisms
\[
  p_i:R(G''_i,E''_i)\xrightarrow{\sim}R.
\]
Thus
\[
  p=p_2^{-1}\circ p_1:
  R(G''_1,E''_1)\xrightarrow{\sim}R(G''_2,E''_2).
\]
The uniqueness theorem gives a unique isomorphism
\[
  f:(G''_1,E''_1)\xrightarrow{\sim}(G''_2,E''_2)
\]
with \(R(f)=p\). This is a descent datum on \(G'\). Its cocycle
condition may be checked after applying \(R\), since \(R\) is fully
faithful, and then follows from the construction of \(p\). Because
\(G'\) is affine, the datum is effective. The pinning is stable
under descent, so it descends to a pinned \(S\)-group \((G,E)\) of
type \(R\).

\medskip
\noindent\textit{N.B.}\quad
In the language of fibered categories, the preceding proof becomes
both shorter and more transparent.

\SGARefTarget{sga3:XXIII:5.10}
\SGAStatement{Corollary 5.10}{Let \(S\) be a nonempty scheme and
\(R\) a pinned root datum for which there exists a reductive
\(S\)-group of type \(R\). Then there is a pinned \(S\)-group of
type \(R\), unique up to unique isomorphism.}

\SGARefTarget{sga3:XXIII:5.11}
\SGAStatement{Definition 5.11}{Under these hypotheses, denote the
unique pinned \(S\)-group of type \(R\) by
\(\operatorname{\acute Ep}_S(R)\), its canonical maximal torus by
\(T_S(R)\), its canonical Borel subgroup by \(B_S(R)\), and so on.}

For every morphism \(S'\to S\), with \(S'\ne\varnothing\), one may
identify
\[
  \operatorname{\acute Ep}_{S'}(R)
    =\operatorname{\acute Ep}_S(R)\times_SS'.
\]
If \(\operatorname{\acute Ep}_{\operatorname{Spec}\mathbf Z}(R)\)
exists, write it \(\operatorname{\acute Ep}(R)\); then
\[
  \operatorname{\acute Ep}_S(R)
    =\operatorname{\acute Ep}(R)
      \times_{\operatorname{Spec}\mathbf Z}S.
\]
The group \(\operatorname{\acute Ep}(R)\) is called the
\emph{Chevalley group scheme of type \(R\)}.

\SGARefTarget{sga3:XXIII:5.12}
\SGAParagraph{5.12}
Thus an \(S\)-sheaf in groups \(G\) is a reductive \(S\)-group of
type \(R\) if and only if it is locally isomorphic, for the étale or
fpqc topology, to \(\operatorname{\acute Ep}_S(R)\). Similarly, by
the conjugacy theorems, giving a reductive \(S\)-group of type \(R\)
with a maximal torus is locally equivalent to
\[
  (\operatorname{\acute Ep}_S(R),T_S(R)).
\]
The analogous statements hold for Borel subgroups and Killing pairs.
